\documentclass[a4paper,fleqn]{cas-sc}

\usepackage{graphicx}
\usepackage{url}
\usepackage[numbers,sort&compress]{natbib}

\begin{document}

\shorttitle{Power Peak Minimization in SALBP}
\shortauthors{Team SALB3PM}

\title[mode = title]{Power Peak Minimization in the Simple Assembly Line Balancing Problem: SAT and MaxSAT Models}

% TODO: Replace the sample author names, emails, and affiliation with the group details.
\author[1]{Team Member 1}
\ead{member1@example.com}
\author[1]{Team Member 2}
\ead{member2@example.com}
% Uncomment and complete this author block if the group has three members.
% \author[1]{Team Member 3}
% \ead{member3@example.com}

\affiliation[1]{organization={Institution / Department},
            city={City},
            country={Country}}

\begin{abstract}
% TODO (about 150--250 words): summarize the SALB3PM problem, implemented
% models, benchmark protocol, principal measured results, and limitations.
This manuscript presents a study of power peak minimization in the simple assembly line balancing problem (SALB3PM). The abstract will be completed after the models and experiments have been implemented and independently validated.
\end{abstract}

\begin{keywords}
Simple assembly line balancing \sep Power peak minimization \sep SAT \sep MaxSAT \sep Pseudo-Boolean encoding
\end{keywords}

\maketitle

\section{Introduction}
% Motivate energy-aware line balancing and state the project contributions.

\section{Problem Definition}
% Define tasks, stations, cycle time, precedence, durations, power, and objective.
% State assumptions and the time discretization used by the implementation.

\section{Related Work}
% Discuss the assigned paper and exact approaches used as baselines.
The primary reference for the MaxSAT formulations studied in this project is the paper by Zheng, Cherif, and Shibasaki~\citep{zheng2024powerpeaks}.

\section{SAT and MaxSAT Models}
\subsection{Task assignment and scheduling constraints}
% Map variables and hard constraints to the paper's notation and equations.

\subsection{Unary peak representation}
% Describe the inaugural model (paper Section III-A / Figure 3).

\subsection{Binary peak representation}
% Describe the binary model (paper Section III-B / Figure 4).

\subsection{Pseudo-Boolean encoding and implementation}
% Record PySAT PBEnc choices, solver versions, and any deviations from the paper.

\section{Exact Baselines}
% Describe CP-SAT, CPLEX CP, CPLEX MIP, and Gurobi implementations actually run.

\section{Experimental Setup}
\subsection{Instances and benchmark provenance}
\subsection{Hardware and software}
\subsection{Protocol and metrics}
% State limits, seeds, repeats, statuses, and validation procedure.

\section{Results}
% Insert tables and figures generated from committed/reproducible raw results.

\section{Discussion and Limitations}
% Interpret only supported observations; discuss license/system limitations.

\section{Conclusion}
% Summarize completed work and future work.

\section*{CRediT authorship contribution statement}
% TODO: Replace with each member's actual contributions.
% Example format: Member Name: Conceptualization, Software, Writing -- original draft.

\section*{Declaration of competing interest}
The authors declare that they have no known competing financial interests or personal relationships that could have appeared to influence the work reported in this paper.

\section*{Data availability}
The source code, experiment scripts, and data supporting this study will be available at \url{<REPOSITORY_URL>}.

\bibliographystyle{cas-model2-names}
\bibliography{references}

\end{document}
